\subsection{正交投影子与对合。}

在本小节中，我们总假设数量 2 在 A 中是 \(invertible\) 的（例如 A 是特征 \(\neq 2\) 的体），并且 \(\Phi\) 是 E 上的一个 \(nondegenerate\) 埃尔米特型。以 \(\frac{1}{2}\) 表示 2 的逆。

引理 2. --- 要使 E 的自同态 u 满足 \(u^{2}=1\) ，必须且只需 \(\frac{1}{2}(1-u)\) 是 E 中的一个投影子；此时 u 是两个投影子 \(\frac{1}{2}(1+u)\) 与 \(\frac{1}{2}(1-u)\) 之差。

实际上，在环 \(\mathfrak{L}(\mathrm{E})\) 中，关系 \(\left(\frac{1}{2}(1 - u)\right)^2 = \frac{1}{2}(1 - u)\) 等价于 \(u^2 = 1\) 。其余是平凡的。

E 的一个自同态 \(u\) （满足 \(u^2 = 1\) ；此时它必然是 E 的等于自身逆的自同构）称为一个对合。令 \(v = \frac{1}{2}(1 - u)\) ，\(U^- = v(E)\) ，\(U^+ = v^{-1}(0)\) （即 \(= w(E)\) ，如果令 \(w = \frac{1}{2}(1 + u)\) ）；我们知道 E 是 \(U^+\) 与 \(U^-\) 的直和（第 VIII 章，§ 1，第 1 小节），并且有 \(u(x) = x\) （在 \(U^+\) 中），\(u(x) = -x\) （在 \(U^-\) 中）。当 A 是一个体且 E 是有限维时，由于 A 的特征 \(\neq 2\) ，由此得出仅有的特征向量 \(\neq 0\) （\(u\) 的）是 U⁺ 或 U⁻ 中的元素 ≠ 0；它们分别对应于特征值 +1 与 −1。

命题 4. --- 设 \(u \in \mathbf{GL}(\mathrm{E})\) 是一个对合。下列性质等价：

\begin{enumerate}
\def\labelenumi{\alph{enumi})}
\item
  u 属于与 \(\Phi\) 相伴的酉群；
\item
  子模 \(U^{+} = \frac{1}{2}(1 + u)(E)\) 与 \(U^{-} = \frac{1}{2}(1 - u)(E)\) 是正交的（从而是非迷向的）。
\end{enumerate}

此外，若 A 是一个体且 E 是有限维的，则性质 a) 与 b) 等价于：

\begin{enumerate}
\def\labelenumi{\alph{enumi})}
\setcounter{enumi}{2}
\tightlist
\item
  \(u = u^{*}\) 。
\end{enumerate}

实际上，对 \(x \in \mathrm{U}^{+}\) 与 \(y \in \mathrm{U}^{-}\) ，关系 \(\Phi(u(x), u(y)) = \Phi(x, y)\) 给出 \(2\Phi(x, y) = 0\) ，因而 \(a)\) 蕴含 \(b)\) 。反之，我们显然有 \(\Phi(u(x), u(y)) = \Phi(x, y)\) （当 \(x\) 与 \(y\) 都属于 \(\mathrm{U}^{+}\) 或都属于 \(\mathrm{U}^{-}\) 时），并且在给定 \(b)\) 时，当二者之一属于 \(\mathrm{U}^{+}\) 而另一个属于 \(\mathrm{U}^{-}\) 时这个关系仍然成立；由于 E 是 \(\mathrm{U}^{+}\) 与 \(\mathrm{U}^{-}\) 的直和，我们看出 \(b)\) 蕴含 \(a)\) 。最后，当 E 是有限维向量空间时，伴随 \(u^{*}\) （\(\Phi\) 非退化时）有定义；关系 \(a)\) 等价于 \(uu^{*} = 1\) ( \(§ 1, n^{\circ} 8\) ，命题 8 的推论)；由于依假设 \(u^{2} = 1\) ，\(a)\) 与 \(c)\) 等价。

推论 1. --- 假设 A 是一个体且 E 是有限维的。映射 \(u \to \frac{1}{2}(1 + u)(E)\) 是从属于与 \(\Phi\) 相伴的酉群的对合 u 所成的集到 E 的非迷向子空间所成的集上的一个双射；与 u 对应的子空间 U \(^+\) 是 E 中在 u 下不变的元素所成的集。

由命题 4，只需证明 E 的每个非迷向子空间 M 都是在某个对合 \(u \in \mathbf{U}(\Phi)\) 下不变的向量所成的集，并且这个对合是唯一的。现在 E 是 M 与 \(\mathbf{M}^0\) 的直和（ \(\S 4\) ，第 1 小节，命题 1 的推论），并且依命题 4，必然有 \(u(x) = x\) （对 \(x \in \mathbf{M}\) ）且 \(u(x) = -x\) （对 \(x \in \mathbf{M}^0\) ）；这些关系唯一地确定 \(u\) ，而且如此确定的自同态 \(u\) 显然满足要求（命题 4）。

我们称如此确定的对合 \(u\) 是关于非迷向子空间 M 的对称。

推论 2. --- 假设 A 是一个体且 E 是有限维的。要使 E 中的投影子 v 满足 v(E) 与 \(v^{-1}(0)\) 正交（从而非迷向），必须且只需 v = v*。

只需对对合 \(u = 1 - 2v\) 应用命题 4。

满足推论 2 的条件的投影子称为对 \(\Phi\) 的正交投影子。

